\subsection{UAV-s-curve}\label{sec:res:uav:scurve}

UAV-s-curve is reported as a caveat and not as a third benchmark. It was piloted with each model at the
budget that model is strongest at --- the rows are not a model comparison --- and every pilot fails --- before any question about the constraint can
be asked of it. A divergence guard ends a flight when the vehicle inverts, and on the plans the models
generate, executed without projection, it ends four to ten flights of ten (\autoref{tab:uav-scurve});
success with constraint satisfaction is $0.00$ in every cell the scene has, projected and unprojected
alike. Because the unprojected plan already fails at the turns at the model's best budget, no smaller
budget was flown and the projected pilots add nothing to the reading: what fails on this scene is the
vehicle's attitude in the two sharp turns, and \autoref{sec:res:uav:controller} isolates the cause by
keeping the plan and changing the tracking controller.

\begin{table}[htpb]
  \caption[UAV-s-curve: the pilot flights]{UAV-s-curve, the pilot flights: each model at the
  budget it is strongest at, without projection, brake-to-rest tracking, ten flights per cell, one
  training seed. \emph{Success}: crosses the finish line; \emph{distance}: mean final distance to the
  goal point of the route; \emph{safe}: no contact and no loss of altitude; \emph{aborted}: ended by the
  divergence guard because the vehicle inverted. \baselinemark\ the baseline of record at its own
  budget.}
  \label{tab:uav-scurve}
  \centering
  \small
  \begin{tabularx}{\linewidth}{@{}L r r r r r r@{}}
    \toprule
    Model & $\nfe$ & Success & Distance [m] & Safe & Aborted & ms/step \\
    \midrule
    MeanFM        & 10 & 0/10 & 2.809 & 0/10 & 10/10 & 88.8 \\
    CI-MeanFM, $\alpha_{\mathrm{end}}=0.2$ & 5 & 2/10 & 2.247 & 2/10 & 8/10 & 45.0 \\
    FM            & 20 & 6/10 & 1.193 & 6/10 & 4/10 & 170.4 \\
    \midrule
    \mainconfig{\baselinemark\ Diffusion} & \mainconfig{20} & \mainconfig{0/10} & \mainconfig{2.432} & \mainconfig{0/10} & \mainconfig{8/10} & \mainconfig{173.9} \\
    \bottomrule
  \end{tabularx}
\end{table}
\dataref{batch\_uav\_20260919\_111701 (tag \texttt{u18sc}: MeanFM $\nfe=10$, FM $\nfe=20$, CI-MeanFM
$\nfe=5$) and batch\_uav\_20260915\_100816 (tag \texttt{u7hg}: diffusion), \texttt{per\_rollout\_detail},
geometry \texttt{s\_curve\_hg}, variant \texttt{diffuser}; \emph{success} = \texttt{success\_relaxed},
\emph{safe} = \texttt{phys\_safe}, \emph{aborted} = \texttt{divergence\_aborted}. These four cells were
Tables 6.14 and C.x of v3.67; the lower budgets that were flown (MeanFM 2, CI-MeanFM 1 and 2, FM 2) are in the
corpus and are no longer printed (author, v3.68). Recounted 2026-09-22.}
\guard{Instantaneous-velocity matching succeeds on six flights of ten, so the controller's limit is not
absolute: a plan whose turns it can follow exists. The two average-velocity models and the baseline invert
on eight to ten flights. Nothing here ranks the models: the budgets differ by row, and what separates the
rows is how the tracker copes with each model's turn, which \autoref{sec:res:uav:controller} examines on
one model with the controller changed.}
\guard{The projected pilots of this scene, flown on the plain constraint set, were withdrawn at v3.67 and
are not re-run: they are not needed to state the caveat. The one key datum the caveat still lacks is the
same unprojected plan under the two other velocity-setpoint policies of the tracker, which is R40 of
\texttt{data\_status/PENDING\_20260923\_uav\_corridor\_v3\_pillars\_v2\_paper\_runs.md}.}

On UAV-s-curve, neither method yields a flight that succeeds within the constraints, so the scene
does not rank them (\autoref{sec:res:uav:scurve}).

\subsection{Caveat: The Tracking Controller}\label{sec:res:uav:controller}

On the manipulator, inverse kinematics and the joint-space controller lie between a generated plan and
the executed motion. The quadrotor is underactuated and
open-loop unstable (\autoref{sec:bg:quadrotor}), and a plan reaches the rotors only through the cascaded
geometric controller that tracks it. That controller needs a velocity setpoint the plan does not supply,
and every result in this chapter uses the brake-to-rest policy of \autoref{sec:method:geometric}: the setpoint
is zero, so the vehicle accelerates towards each new position and brakes to rest at it. The result of a flight is therefore a property of the planner and the controller
together, and UAV-s-curve, with its two sharp turns, is where the controller becomes the limit.

This is shown by keeping the plan and changing the controller (\autoref{tab:uav-controller}). Analytic
average-velocity matching at $\nfe=10$, executed without projection, succeeds on none of ten flights
with the brake-to-rest controller and ends on average $2.81$\,m from the goal point of its route. With
MuJoCo MPC \parencite{howell2022predictive} as the tracking controller, the same plans end $0.299$, $0.298$
and $0.294$\,m from that point, and two of the three flights cross the finish line. The brake-to-rest
failures are not wall collisions: on all ten flights the vehicle stays at altitude until it loses attitude
control and inverts, between control steps $395$ and $421$, about $12$\,s into the flight. MuJoCo MPC flies
the same turn. The unprojected plan is flyable; the brake-to-rest controller cannot fly it.

What separates the two is the velocity setpoint and not the gains. Both controllers act on the same
vehicle with the same nominal gains, which the evaluation fixes and never varies, and the brake-to-rest
policy is the one that tracks the reference most closely: its mean tracking error, $0.29$ to $0.34$\,m,
is the smallest of any configuration on this scene, and it is also the only one that never advances along
the route. A controller that follows a reference tightly while making no progress along it is not failing
to track; it is failing to move. On a scene whose reference demands sustained forward speed, a setpoint
of zero velocity at every plan step is what produces both that standstill and the loss of attitude that
ends the flight.

\begin{table}[htpb]
  \caption[UAV-s-curve: the same plan under two tracking controllers]{UAV-s-curve, analytic
  average-velocity matching at $\nfe=10$, the same plans under two tracking controllers.
  \emph{Cascaded geometric, brake-to-rest} is the controller behind every other result in this chapter;
  \emph{MuJoCo MPC} (\autoref{sec:method:mjpc}) replaces it here and nowhere else. The last column is
  the measured elapsed time per executed control step, including planner and controller. \emph{Success}: crosses
  the finish line; \emph{distance}: mean final distance to the goal point of the route; \emph{safe}: no
  contact and no loss of altitude. One training seed.}
  \label{tab:uav-controller}
  \centering
  \small
  \begin{tabularx}{\linewidth}{@{}L L r r r r r@{}}
    \toprule
    Projection & Controller & Flights & Success & Distance [m] & Safe & ms/step \\
    \midrule
    none                   & cascaded geometric, brake-to-rest & 10 & 0/10 & 2.809 & 0/10 & 94.9 \\
    none                   & MuJoCo MPC    & 3  & 2/3  & 0.297 & 2/3 & 215.6 \\
    per-step, best rule    & cascaded geometric, brake-to-rest & 10 & 0/10 & 2.446 & --- & --- \\
    per-step, random       & MuJoCo MPC    & 3  & 0/3  & 2.571 & --- & 2660.1 \\
    \bottomrule
  \end{tabularx}
\end{table}
\dataref{uav\_results.py, s-curve block; Campaign DA\_20260909\_T5\_mjpc\_vs\_pid\_s\_curve F1--F3, F7;
Report\_20260909\_MJPC\_vs\_PID\_s\_curve \S5.4: full-step times on the unprojected
arms and the MuJoCo MPC per-step random arm; the other projected controller cell has no matched
elapsed-time figure in this report}
\guard{Three flights with MuJoCo MPC. The gains are held at their nominal values in every evaluation of
this thesis and were not tuned; the two alternative settings that exist scale the position loop by
$\pm20\,\%$ and leave the attitude loop and the thrust limit untouched, so neither can be expected to
address a loss of attitude. What this section supports is that the scene is limited by the choice of
velocity setpoint, not that no tuning of this controller could fly it.}

\begin{figure}[htpb]
  \centering
  \includegraphics[width=\linewidth]{fig_uav_scurve_paths}
  \caption[UAV-s-curve: flown paths under the two tracking controllers]{UAV-s-curve from above, drawn like
  \autoref{fig:constraints-uav}, with the executed path of every flight of the two unprojected rows of
  \autoref{tab:uav-controller}: the same plans under the two tracking controllers. Green: success; red:
  failure; a cross at the end of a path marks a flight that crashed or that the divergence guard ended.}
  \label{fig:uav-scurve-paths}
\end{figure}
\dataref{DA\_in\_Paper/data/uav\_paths.json, extract/uav\_paths.py; runs
Emf\_K10\_mpc4\_mjpc\_T0.5\_u7hg and its \texttt{pid\_stopgo} sibling}

In that figure a brake-to-rest path stops where the vehicle lost its attitude rather than where it ran
out of authority to track: all ten flights end at the guard, while MuJoCo MPC records no abort,
follows the second passage to the goal on two flights of three, and on the third reaches the finish
line's $x$ with a contact. Neither controller is collision-free
without projection, which is what the projected rows of \autoref{tab:uav-controller} address.

Changing the controller is not sufficient once projection changes the plan. With per-step projection the
plan succeeds on none of three flights under MuJoCo MPC and on none of ten under the brake-to-rest
controller, and no flight succeeds within the constraints. The unprojected experiment isolates a limit of
the controller; the projected experiment remains a failure of the planner, the projection and the
controller together. UAV-s-curve therefore supports no ordering of the generative models and no preferred
projection method, and the ordering on the two other scenes is to be read for the controller with which it
was measured.

The other reason MuJoCo MPC replaces the controller here and nowhere else is what it costs to run. The
time reported in every other table of this chapter is measured around the planner and stops there
(\autoref{sec:setup:metrics:avoiding}), so it prices no controller at all; pricing one means reading the
elapsed time of the flights themselves and dividing it by the control steps that executed
(\autoref{tab:uav-controller-cost}). On the unprojected arm, where nothing lies between the plan and the
controller, a control step takes $94.9$\,ms under the brake-to-rest controller and $215.6$\,ms under
MuJoCo MPC. The planner's share of that is the same in both, between $88$ and $90$\,ms at $\nfe=10$, so
the difference belongs to the controller: about $6$\,ms per step for the cascaded geometric law against
about $126$\,ms for MuJoCo MPC, which solves an optimisation problem at every control step where the
other evaluates an algebraic law. The controller that flies this scene is about twenty times the price
of the one that does not, and by itself it costs more per control step than generating the plan --- some
seven times what the operating point of D3IL-avoiding spends on an entire action
(\autoref{tab:avoiding-conclusion}).

\begin{table}[htpb]
  \caption[UAV-s-curve: what each tracking controller costs per control step]{UAV-s-curve, MeanFM at $\nfe=10$
  without projection: the elapsed time of a flight divided by the control steps it executed, in
  milliseconds, under the two tracking controllers of \autoref{tab:uav-controller} --- the one quantity
  in this chapter that contains the controller. \emph{The planner}: the planner time reported elsewhere;
  \emph{the controller and the simulator step}: the remainder. Ten flights on the left, three on the
  right, one training seed.}
  \label{tab:uav-controller-cost}
  \centering
  \small
  \begin{tabularx}{\linewidth}{@{}L r r@{}}
    \toprule
    Time per executed control step [ms] & Cascaded geometric, brake-to-rest & MuJoCo MPC \\
    \midrule
    The control loop                          & $94.9$ & $215.6$ \\
    \quad the planner                         & $88.4$ & $89.8$ \\
    \quad the controller and the simulator step & $6.5$ & $125.8$ \\
    \bottomrule
  \end{tabularx}
\end{table}
\dataref{Report\_20260909\_MJPC\_vs\_PID\_s\_curve \S5.4--5.5, from the job records of Campaign
DA\_20260909\_T5\_mjpc\_vs\_pid\_s\_curve; planner time from the same batch as
\autoref{tab:uav-controller}}
\guard{This is a decomposition of recorded elapsed times and not a measurement taken around the
controller call, which the evaluation does not instrument. The remainder row therefore carries the
simulator step with the controller, and it is only their \emph{difference} that isolates the controller.
One training seed, ten flights against three.}

What that price buys is a traverse and not a result. Neither controller produces a flight that succeeds
within the constraints on this scene: MuJoCo MPC turns flights that never arrive into flights
that arrive while still violating, at about twenty times the cost per control step of the controller
used everywhere else in this thesis. It is used here to locate the limit, and it is not adopted.
